\begin{center}\bfseries 24.\end{center}
\noindent\hspace*{\parindent}Von dem, was \textsc{Gauss} in der Theorie der kubischen Reste, die er von Anfang an zugleich mit derjenigen der biquadratischen durchdachte, erarbeitet hat, geben uns nur die zuvor erwähnten Fragmente (W. VIII, S. 5—14) einige Kunde. Von ihnen finden sich die auf S. 5—8 abgedruckten in einem mit der Überschrift \textit{Uraniae sacrum} versehenen, im November 1802 begonnenen Hefte (Al) des Nachlasses unmittelbar auf astronomische und magnetische Rechnungen folgend, die bis 1805 ausgedehnte Beobachtungen benutzen; sie werden also vermutlich nicht sehr viel später aufgezeichnet sein. Teils hat \textsc{Gauss} in ihnen den kubischen Charakter einer Zahl (mod. \(p\)) in Zusammenhang gebracht mit der Darstellbarkeit von \(p\) durch gewisse binäre quadratische Formen, teils ihn mit Hilfe der Kreisteilung festzustellen gelehrt. So wird neben einem ähnlichen Satze über biquadratische Reste folgende Aussage von ihm gegeben (W. VIII, S. 7): Seien \(p\equiv1,\ q\equiv\pm1\) (mod. 3) zwei verschiedene Primzahlen und es werde entsprechend den doppelten Vorzeichen
\[
(a+3b\sqrt{-3})^{\frac{q\mp1}{3}}\equiv A\pm B\sqrt{-3}\ (\mathrm{mod.}\ q)
\]
gesetzt; je nach den drei möglichen Fällen
\[
B\equiv0,\ B\equiv A,\ B\equiv-A\ (\mathrm{mod.}\ q)
\]
ist dann
\[
q^{\frac{p-1}{3}}\equiv1,\ -\frac12-\frac12\mathbin{.}\frac{a}{3b},\ -\frac12+\frac12\mathbin{.}\frac{a}{3b}\ (\mathrm{mod.}\ p).
\]

In dieser Aussage liegt unmittelbar, wenn man sich \(p\) in seine komplexen Faktoren zerlegt denkt, das kubische Fundamentaltheorem für den Fall
\newpage\SourceLeaf{67}\ArithmeticHead{60}{L}{BACHMANN, ÜBER GAUSS’ ZAHLENTHEORETISCHE ARBEITEN.}
\noindent \(q\equiv-1\) (mod. 3). \textsc{Gauss}’ im Nachlassteile Ec 1 befindlicher, W. VIII, S. 9 abgedruckter Beweis desselben aber, der später gefunden sein wird als die vorerwähnten Untersuchungen, entspringt aus dem gleichen Ideenkreise, wie sein Beweis VI des quadratischen und der erwähnte Beweis des biquadratischen Gesetzes, nämlich aus den Eigenschaften des Ausdrucks
\[
r+\rho r^g+\rho^2r^{g^2}+\ldots+\rho^{p-2}r^{g^{p-2}},
\]
wo \(g\) eine primitive Wurzel (mod. \(p\)) und \(\rho\) eine kubische Einheitswurzel bedeutet. Da nach den \textsc{Fricke}schen Erläuterungen (W. VIII, S. 11) aus jener Aussage auch für den Fall \(q\equiv+1\) (mod. 3) das kubische Reziprozitätsgesetz, wie es später von \textsc{Jacobi}\footnote{\textsc{Crelles} Journal für Mathematik \textbf{2} (1827), S. 66, Werke VI, S. 233.}\ und \textsc{Eisenstein}\begingroup\renewcommand{\thefootnote}{**)}\footnote{\textsc{Crelles} Journal für Mathematik \textbf{27}, S. 289; \textbf{28}, S. 28 (1844).}\endgroup\ bewiesen wurde, sich erschliessen lässt, ist mittelbar durch die \textsc{Gauss}sche Methode dieses Gesetz schon vordem in seiner ganzen Allgemeinheit festgestellt worden.

Ein weiteres Fragment (abgedruckt W. VIII, S. 15) zeigt \textsc{Gauss} bemüht, die Methode seiner Beweise III und V durch geeignete Einteilung des gesamten Restsystems eines Moduls in Drittel, wobei die Gittervorstellung benutzt wird, auf kubische Reste auszudehnen. Dass ein auf die Zahl Drei und einen Modul \(p=3n+1\) bezüglicher Beweis ihm in ähnlicher Weise geglückt sein muss, zeigt die vom 6. Januar 1809 datierte Tagebuchnotiz (Nr. 138 T.), aus der wohl zu schliessen ist, dass das letztgenannte Fragment vor diesem Datum verfasst worden ist. Andererseits kann es nicht vor dem August 1808 geschrieben sein, da es in \textsc{Gauss}’ Handbuch 18, Bd, vom Oktober 1805 erst S. 205—213 auf astronomische Beobachtungen aus diesem Monate folgt.
